\documentclass[11pt]{article}
\usepackage[margin=1in]{geometry}
\usepackage{graphicx}
\usepackage{amsmath, amssymb}
\usepackage{booktabs}
\usepackage{microtype}
\usepackage{caption}
\usepackage{xcolor}
\usepackage{hyperref}
\hypersetup{colorlinks=true, linkcolor=blue!50!black,
            urlcolor=blue!50!black, citecolor=blue!50!black}

\graphicspath{{figures/}}

\newcommand{\figfit}[1]{\includegraphics[width=\linewidth,
   height=0.78\textheight,keepaspectratio]{#1}}

\title{Auditory superstitious-perception task:\\
       a focused report on Full-vs-Imagined block, site, and block-order
       effects}
\author{}
\date{}

\begin{document}
\maketitle

\begin{abstract}
\noindent
This report examines the auditory superstitious-perception dataset
($N\!=\!26$ subjects: 16 Oberlin, 10 UChicago, two 150-trial blocks each)
and focuses on 3 main quesitons: (i) does behaviour differ
between the \texttt{full\_sentence} and \texttt{imagined\_sentence}
blocks; (ii) do the two recruitment sites differ; and (iii) does
block-order matter.  Every analysis is repeated three times
(UChicago alone, Oberlin alone, and unified pool) so that any
interpretation can be cross-checked across experimental sites.
The headline result is that \emph{the Oberlin full-sentence block is
significant at the group level and this is verified multiple statistical tests}; 
UChicago is null; the site difference is trending in the
expected direction and a sub-sampling power analysis shows this
trend is consistent with a real but under-powered effect.  Block
order shows no reliable effects.  Questionnaire correlations are all
exploratory and none survives correction.
\end{abstract}

\tableofcontents

%==============================================================================
\section{Reading guide}
%==============================================================================

The task asks each subject to indicate, on each trial, whether a white noise
auditory stimulus was the experimenter-defined ``target'' or ``distractor''. The two block types are
\texttt{full\_sentence} and \texttt{imagined\_sentence}.  
Subjects also completed eight personality / state
questionnaires.

\paragraph{Three-way reporting convention.}
Every comparison in this report is run three times.
\begin{itemize}
\item \textbf{UChicago} alone ($n=10$),
\item \textbf{Oberlin}  alone ($n=16$),
\item \textbf{Unified}  pool ($n=26$).
\end{itemize}
Where a regression line is drawn on a scatter plot, each of the three
groups gets its \emph{own panel} with its \emph{own scatter} and its
\emph{own ordinary-least-squares line}.  We never overlay groups onto
one panel with a single shared line, because that obscures whether the
relationship is consistent across sites or driven by one of them.

\paragraph{Statistical conventions.}
For each comparison we report:
\begin{itemize}
\item the sample size,
\item the observed mean (and 95\% bootstrap CI where natural),
\item a non-parametric test (one-sample Wilcoxon vs.\ a null value, or
      Mann--Whitney U for unpaired groups, or paired Wilcoxon for
      within-subject contrasts),
\item the parametric counterpart ($t$-test) for context,
\item Cohen's $d$ as a unit-free effect size,
\item where relevant, a permutation $p$-value with $10^4$ resamples.
\end{itemize}
We use Wilcoxon / Mann--Whitney as the primary tests because $N$ is
small and the distributions are not guaranteed normal.

\paragraph{What ``Cohen's $d$'' means in plain words.}
Cohen's $d$ is the size of an effect measured in standard-deviation
units.  $|d|=0.2$ is small, $|d|=0.5$ is medium, $|d|=0.8$ is large [cohen 1988].
An effect of $d=1.0$ is one standard deviation
of the data --- an ``obvious'' difference.  Reporting $d$ alongside
$p$-values is important because $p$ depends heavily on $N$ but $d$
does not; a large $d$ with a marginal $p$ usually means the experiment
is under-powered, not that the effect is absent.

\paragraph{What a non-parametric test is doing.}
The Wilcoxon signed-rank test ranks the $|x_i - \mu_0|$ values and
sums the ranks attached to positive deviations.  Under the null
$x_i \overset{d}{=} 2\mu_0 - x_i$ this rank-sum has a known
distribution, so a $p$-value falls out without assuming Gaussian
data.  The Mann--Whitney $U$ test is its two-sample analogue.  For our
small $N$ they are the right default.

%==============================================================================
\section{Section A.  Accuracy vs chance, per group, per block}
%==============================================================================

This is the base question: in each block, in each group, are subjects
performing reliably above chance ($\text{accuracy} = 0.5$)?

\begin{table}[h]
\centering
\small
\begin{tabular}{llrrrrrrr}
\toprule
block & group & $n$ & mean & 95\% CI & Wilcoxon $p$ & $t$-test $p$ & Cohen's $d$ \\
\midrule
full\_sentence     & UChicago & 10 & 0.504 & [0.479,0.530] & 0.904 & 0.777 & $+0.09$ \\
full\_sentence     & Oberlin  & 16 & 0.526 & [0.513,0.538] & \textbf{0.005} & \textbf{0.001} & $\mathbf{+0.99}$ \\
full\_sentence     & Unified  & 26 & 0.518 & [0.504,0.531] & \textbf{0.026} & \textbf{0.017} & $+0.50$ \\
imagined\_sentence & UChicago & 10 & 0.490 & [0.463,0.515] & 0.611 & 0.487 & $-0.23$ \\
imagined\_sentence & Oberlin  & 16 & 0.510 & [0.495,0.526] & 0.350 & 0.238 & $+0.31$ \\
imagined\_sentence & Unified  & 26 & 0.502 & [0.488,0.516] & 0.620 & 0.758 & $+0.06$ \\
\bottomrule
\end{tabular}
\caption{Section A: one-sample tests of accuracy against chance ($\mu_0 = 0.50$).
Bold rows are significant at $\alpha = 0.05$.  CI is a $10^4$-bootstrap
percentile interval on the mean.  See
\texttt{tables/A\_accuracy\_vs\_chance.csv}.}
\end{table}

\begin{figure}[h]
\centering
\figfit{A_accuracy_vs_chance.png}
\caption{Section A.  Accuracy distributions for each (group, block).
Red dashed line is chance ($0.50$); black solid line is the group mean.
Title shows the bootstrap CI on the mean and the Wilcoxon $p$ vs
chance.}
\end{figure}

\paragraph{Read.}
Two cells are nominally significant at $\alpha=0.05$ without
correction: \textbf{Oberlin full\_sentence} ($p=0.005$) and
\textbf{Unified full\_sentence} ($p=0.026$).  Under Bonferroni
correction for 6 cells (threshold $\alpha = 0.05/6 = 0.0083$), only
\textbf{Oberlin full\_sentence} survives.  Its effect size is large
($d \approx 1.0$), so this is not a fragile $p$-value.  The Unified
pool inherits the result but at a smaller effect size, because the
UChicago subjects are at chance and they dilute the group mean.
Imagined-sentence is at chance for everyone.

%==============================================================================
\section{Section B.  Full sentence vs Imagined sentence (paired)}
%==============================================================================

Each subject ran both blocks, so the right test is paired.  We also
report an unbiased permutation test (sign-flip the within-subject
differences $10^4$ times).

\begin{table}[h]
\centering
\small
\begin{tabular}{lrrrrrrr}
\toprule
group & $n$ & mean(F) & mean(I) & mean diff & paired Wilcoxon $p$ & perm $p$ & $d$ \\
\midrule
UChicago & 10 & 0.504 & 0.490 & $+0.014$ & 0.543 & 0.508 & $+0.22$ \\
Oberlin  & 16 & 0.526 & 0.510 & $+0.016$ & 0.187 & 0.191 & $+0.35$ \\
Unified  & 26 & 0.518 & 0.502 & $+0.015$ & 0.149 & 0.154 & $+0.29$ \\
\bottomrule
\end{tabular}
\caption{Section B: full-sentence minus imagined-sentence, paired.
See \texttt{tables/B\_full\_vs\_imagined.csv}.}
\end{table}

\begin{figure}[h]
\centering
\figfit{B_full_vs_imagined.png}
\caption{Section B.  Per-subject within-subject lines connecting Full
to Imagined accuracy, in each of the three reporting groups.  Coloured
heavy line is the group mean.}
\end{figure}

\paragraph{Read.}
The within-subject difference is $\approx +0.015$ in every group, in
the expected direction (Full $>$ Imagined), but never reaches
$p < 0.05$.  The effect sizes are small-to-medium.  The fact that the
direction is the same in all three reporting groups argues this is a
genuine micro-effect rather than noise; Section G below quantifies how
many more subjects we would need to declare it.

%==============================================================================
\section{Section C.  Oberlin vs UChicago site comparison}
%==============================================================================

Per block, an unpaired comparison.

\begin{table}[h]
\centering
\small
\begin{tabular}{lrrrrrrrr}
\toprule
block & $n_O$ & $n_U$ & mean$_O$ & mean$_U$ & diff & MWU $p$ & Welch $p$ & $d$ \\
\midrule
full\_sentence     & 16 & 10 & 0.526 & 0.504 & $+0.022$ & 0.145 & 0.173 & $+0.65$ \\
imagined\_sentence & 16 & 10 & 0.510 & 0.490 & $+0.020$ & 0.475 & 0.231 & $+0.54$ \\
\bottomrule
\end{tabular}
\caption{Section C: site comparison per block (unpaired).  See
\texttt{tables/C\_site\_comparison.csv}.}
\end{table}

\begin{figure}[h]
\centering
\includegraphics[width=0.9\linewidth]{C_site_comparison.png}
\caption{Section C.  Per-block site boxplots with overlaid jittered
points.}
\end{figure}

\paragraph{Read.}
Oberlin is numerically higher than UChicago in both blocks
($+0.022$ in full, $+0.020$ in imagined) with \emph{medium} Cohen's
$d$ ($\approx 0.55$--$0.65$).  Neither comparison reaches $p < 0.05$.
Whether this is a real site effect that simply needs more subjects, or
just sampling noise, is the question Section G addresses
quantitatively; the short answer is ``consistent with a real
under-powered medium-sized effect''.

%==============================================================================
\section{Section D.  Block order}
%==============================================================================

\begin{table}[h]
\centering
\small
\begin{tabular}{llrrrrrr}
\toprule
block & group & $n_{B1}$ & $n_{B2}$ & mean(B1) & mean(B2) & MWU $p$ & $d$ \\
\midrule
full\_sentence     & UChicago &  4 &  6 & 0.500 & 0.507 & 0.831 & $+0.15$ \\
full\_sentence     & Oberlin  &  5 & 11 & 0.519 & 0.529 & 0.493 & $+0.35$ \\
full\_sentence     & Unified  &  9 & 17 & 0.511 & 0.521 & 0.432 & $+0.29$ \\
imagined\_sentence & UChicago &  6 &  4 & 0.490 & 0.490 & 0.915 & $0.00$ \\
imagined\_sentence & Oberlin  & 11 &  5 & 0.500 & 0.532 & 0.078 & $+1.03$ \\
imagined\_sentence & Unified  & 17 &  9 & 0.497 & 0.513 & 0.185 & $+0.43$ \\
\bottomrule
\end{tabular}
\caption{Section D: accuracy by block-order index, per group.  See
\texttt{tables/D\_block\_order.csv}.}
\end{table}

\paragraph{Read.}
Nothing reaches $\alpha=0.05$.  The strongest single number is
imagined-sentence in Oberlin, where the eleven subjects who ran
imagined first averaged $0.500$ and the five who ran it second
averaged $0.532$ ($p=0.08$, $d=1.03$).  That is the kind of pattern
practice/familiarisation would produce, but with $n_{B2}=5$ this is ``suggestive only''.

%==============================================================================
\section{Section E.  Convergent metrics: Oberlin full\_sentence}
%==============================================================================

The single most important question of the project is whether the
Oberlin full-sentence accuracy result is real or a one-off. The
prior data anlysis computed several
\emph{mathematically distinct} per-subject quantities, all of which
should equal zero if subjects are responding at random and all of
which should be positive if subjects are exploiting the experimenter's
acoustic structure.  A result that shows up in only one of these
metrics is suspect; a result that shows up in all of them is robust.
We re-evaluate four such metrics here, restricted to the
\texttt{full\_sentence} block, with the standard three-way split.

\paragraph{The four metrics, in plain words.}
\begin{itemize}
\item \textbf{$d'$ (signal-detection sensitivity)} measures how
      separable a subject's distributions of internal evidence are on
      true-target vs.\ true-distractor trials.  Formally, $d' =
      \Phi^{-1}(\text{hit rate}) - \Phi^{-1}(\text{false-alarm rate})$.
      $d'=0$ means no discrimination; $d'>0$ means the subject is on
      average more likely to say ``target'' on true targets.
\item \textbf{Accuracy} is just $(\text{hits} + \text{correct
      rejections})/N$.  Mathematically different from $d'$ but
      derived from the same $2\times 2$ table.
\item \textbf{$\cos(\mathrm{CI}_s, \mathrm{CI}_{\rm true})$} is the
      cosine between the subject's classification image (the average
      raw waveform of stimuli they called ``target'' minus the average
      of those they called ``distractor'') and the experimenter's
      target-minus-distractor template.  It lives in 2960-D raw audio
      space; it is positive when the subject's implicit template
      points the same direction as the true template.
\item \textbf{$\cos(w_s, w_{\rm true})$} is the same thing but in a
      32-band log-spectral basis: each subject's ridge-logistic
      weights $w_s$ are compared to the analogous weights trained on
      the experimenter labels.
\item \textbf{$\beta_{\text{truetarget}}$} is the per-subject
      coefficient on the experimenter label in a logistic regression
      of called\_target on temporal covariates plus the true label.
      It is positive when the subject's $P(\text{yes})$ is higher on
      true targets than on true distractors, after partialling out
      reminder-cycle and block-position effects.
\end{itemize}
These five quantities are computed on overlapping data but they are
\emph{not} statistically equivalent.  $d'$ and accuracy come from a
$2\times 2$ contingency table; the two cosines come from
high-dimensional templates that incorporate the acoustic content of
each individual stimulus; $\beta_{\text{truetarget}}$ comes from a
logistic with extra nuisance regressors.  A consistent group result
across all of them is the strongest convergent evidence the dataset
can offer.

\begin{table}[h]
\centering
\small
\begin{tabular}{llrrrrr}
\toprule
group & metric & $n$ & mean & 95\% CI & Wilcoxon $p$ & $d$ \\
\midrule
UChicago & $d'$               & 10 & $+0.021$ & [$-0.10$,$+0.15$] & 0.846 & $+0.09$ \\
UChicago & accuracy           & 10 & $+0.504$ & [$+0.48$,$+0.53$] & 0.904 & $+0.09$ \\
UChicago & $\cos(\rm CI)$     & 10 & $-0.002$ & [$-0.06$,$+0.06$] & 1.000 & $-0.02$ \\
UChicago & $\cos(w)$          & 10 & $+0.036$ & [$-0.07$,$+0.14$] & 0.625 & $+0.20$ \\
UChicago & $\beta_{\rm true}$ & 10 & $+0.056$ & [$-0.13$,$+0.25$] & 0.625 & $+0.17$ \\
\midrule
\textbf{Oberlin} & $d'$               & 16 & $\mathbf{+0.135}$ & [$+0.07$,$+0.20$] & \textbf{0.001} & $\mathbf{+0.99}$ \\
\textbf{Oberlin} & accuracy           & 16 & $\mathbf{+0.526}$ & [$+0.51$,$+0.54$] & \textbf{0.005} & $\mathbf{+0.99}$ \\
\textbf{Oberlin} & $\cos(\rm CI)$     & 16 & $\mathbf{+0.057}$ & [$+0.03$,$+0.08$] & \textbf{0.001} & $\mathbf{+1.02}$ \\
Oberlin          & $\cos(w)$          & 16 & $+0.081$ & [$-0.02$,$+0.18$] & 0.231 & $+0.38$ \\
\textbf{Oberlin} & $\beta_{\rm true}$ & 16 & $\mathbf{+0.218}$ & [$+0.12$,$+0.32$] & \textbf{0.003} & $\mathbf{+1.05}$ \\
\midrule
Unified  & $d'$               & 26 & $+0.091$ & [$+0.02$,$+0.16$] & 0.013 & $+0.51$ \\
Unified  & accuracy           & 26 & $+0.518$ & [$+0.50$,$+0.53$] & 0.026 & $+0.50$ \\
Unified  & $\cos(\rm CI)$     & 26 & $+0.034$ & [$+0.00$,$+0.06$] & 0.025 & $+0.44$ \\
Unified  & $\cos(w)$          & 26 & $+0.064$ & [$-0.01$,$+0.14$] & 0.150 & $+0.32$ \\
Unified  & $\beta_{\rm true}$ & 26 & $+0.156$ & [$+0.05$,$+0.25$] & 0.011 & $+0.59$ \\
\bottomrule
\end{tabular}
\caption{Section E: four convergent alignment statistics in the
\texttt{full\_sentence} block, three-way split.  Bold = significant at
$\alpha=0.05$ for that group--metric combination.  All five Oberlin
means are positive; four of five are individually significant.
\texttt{tables/E\_convergent\_full\_sentence.csv}.}
\end{table}

\begin{figure}[h]
\centering
\figfit{E_convergent_full_sentence.png}
\caption{Section E.  Four convergent metrics for the full-sentence
block.  Rows: metric.  Columns: UChicago, Oberlin, Unified.  Red
dashed line is the null value (0 for $d'$ and the cosines, 0.5 for
accuracy); black line is the group mean.  Each panel reports its own
mean, Wilcoxon $p$, and Cohen's $d$.}
\end{figure}

\paragraph{Read.}
For Oberlin in the full-sentence block, four mathematically distinct
quantities ($d'$, accuracy, $\cos(\mathrm{CI})$,
$\beta_{\text{truetarget}}$) all yield $p<0.01$ with $d \approx 1.0$;
the fifth ($\cos(w)$ in the spectral basis) points the same way with
$d=0.38$ but does not reach significance with $n=16$.  This is the
kind of \emph{mutually corroborating} pattern that one rarely gets
from noise: a noisy null result will generate one significant cell
out of five $\sim 23\%$ of the time, but it will essentially never
generate four significant cells with all five effect signs the same.
For UChicago, in contrast, all five metrics are at or near zero with
small effect sizes; the Unified pool inherits a diluted version of
the Oberlin result.

%==============================================================================
\section{Section F.  Questionnaire correlations (exploratory)}
%==============================================================================

For each questionnaire we compute Spearman $\rho$ with accuracy in
each group and each block.  This is exploratory: there are
$8\times 2\times 3 = 48$ such correlations, and Bonferroni at $\alpha
= 0.05$ would require $p < 1.04\times 10^{-3}$.  No correlation
clears this bar.  We report the strongest in the table, then plot
\emph{every} questionnaire as scatter+OLS line in three side-by-side
panels (UChicago, Oberlin, Unified), one full figure per block.

\begin{table}[h]
\centering
\small
\begin{tabular}{lllrrr}
\toprule
block & group & quest & $n$ & $\rho$ & $p$ \\
\midrule
full\_sentence     & Unified  & BAIS\_C & 26 & $+0.31$ & 0.122 \\
full\_sentence     & UChicago & LSHS    & 10 & $+0.52$ & 0.124 \\
full\_sentence     & Oberlin  & FSS     & 16 & $+0.37$ & 0.156 \\
full\_sentence     & UChicago & BAIS\_C & 10 & $+0.48$ & 0.160 \\
full\_sentence     & Unified  & BAIS\_V & 26 & $+0.28$ & 0.169 \\
full\_sentence     & Unified  & FSS     & 26 & $+0.28$ & 0.170 \\
full\_sentence     & Oberlin  & BAIS\_C & 16 & $+0.36$ & 0.177 \\
\bottomrule
\end{tabular}
\caption{Section F: top eight Spearman $\rho$ between a questionnaire
and accuracy.  None survives correction.  Full table:
\texttt{tables/F\_questionnaire\_spearman.csv}.}
\end{table}

\begin{figure}[p]
\centering
\figfit{F_quest_vs_accuracy_full_sentence.png}
\caption{Section F.  Each panel: one group, one questionnaire, scatter
of (questionnaire, accuracy) plus a panel-specific OLS line.  Red
dotted line is chance accuracy.}
\end{figure}

\begin{figure}[p]
\centering
\figfit{F_quest_vs_accuracy_imagined_sentence.png}
\caption{Section F (continued).  Same layout for the
imagined-sentence block.}
\end{figure}

\paragraph{Read.}
The cleanest \emph{descriptive} pattern is BAIS\_C
(Bucknell Auditory Imagery -- Control) showing a positive trend in
all three groups for the full-sentence block ($\rho = +0.48$ in
UChicago, $+0.36$ in Oberlin, $+0.31$ in Unified).  None of these
are individually significant.  Book 3 of this project (Ch.\ 28-29)
reported a separate, complementary pattern in BAIS\_V where 7 of 8
permutation-significant subject-blocks scored above the median on
auditory-imagery vividness.  Both observations are suggestive of an
auditory-imagery component to who-can-do-this-task; neither is
statistically confirmed.

%==============================================================================
\section{Section G.  Is it a power problem?  Sub-sample diagnostics}
%==============================================================================

Several of the comparisons above are ``trending'' --- the direction
is right, the effect size is medium-to-large, but the $p$-value is
in $[0.05, 0.20]$.  The right diagnostic question is: \emph{at the
observed effect size, would more subjects have made this significant,
or is the effect itself too small to be reliable?}  We answer by
sub-sampling.  At each candidate sub-sample size $n$, we draw 2000
random subsets of that size from the observed data, run the relevant
test, and record the fraction with $p < 0.05$.  This estimates the
power that the observed effect would deliver at sample size $n$.

\paragraph{Three diagnostics.}
\begin{itemize}
\item \textbf{G1.} Oberlin full\_sentence accuracy vs.\ chance.
      Observed $N=16$, $p=0.005$, $d=0.99$.  We expect this to be
      easily powered.
\item \textbf{G2.} Site difference (Oberlin vs UChicago),
      full\_sentence accuracy.  Observed total $N=26$, $p=0.145$,
      $d=0.65$.
\item \textbf{G3.} Full vs Imagined paired (Unified).  Observed
      $N=26$, $p=0.149$, $d=0.29$.
\end{itemize}

\begin{figure}[h]
\centering
\figfit{G_power_curves.png}
\caption{Section G.  Sub-sample power curves.  In each panel the
$x$-axis is the sub-sample size and the $y$-axis is the empirical
probability that a re-test on that subsample would yield $p<0.05$.
Red dashed line is the conventional $0.80$ power threshold.}
\end{figure}

\begin{table}[h]
\centering
\small
\begin{tabular}{lrrrr}
\toprule
comparison & observed $n$ & observed $p$ & observed $d$ & $n$ for 80\% power \\
\midrule
G1: Oberlin full\_sentence vs chance     & 16 & 0.005 & $+0.99$ & 11 \\
G2: Site Oberlin vs UChicago, full       & 26 & 0.145 & $+0.65$ & not reached at $N=26$ \\
G3: Full vs Imagined paired (Unified)    & 26 & 0.149 & $+0.29$ & not reached at $N=26$ \\
\bottomrule
\end{tabular}
\caption{Section G: empirical power summary.  See
\texttt{tables/G\_required\_n\_summary.csv}.}
\end{table}

\paragraph{Read.}
G1 confirms that Oberlin-full-sentence effect is well-powered and even an $n=11$ subsample would, on average,
detect it.

For both G2 and G3 the observed effect size is non-trivial ($d=0.65$ for the
site comparison; $d=0.29$ for the paired Full vs.\ Imagined
contrast) and the observed direction is the predicted one, but the
empirical power at $N=26$ is below $0.80$.  In other words, the
non-significant $p$-values are exactly what we should \emph{expect}
to see for a real effect of this magnitude with this sample size.
Power calculations confirm that, holding $d=0.65$,
roughly $n_O = n_U \approx 37$ each would be needed to push a
two-sample test to 80\% power; for the within-subject
$d=0.29$ contrast, roughly $N \approx 93$. (DOUBLE CHECK URSELF HERE)

\paragraph{Honest qualifier.}
Empirical sub-sampling necessarily holds the observed effect size
fixed.  If the true effect is smaller than what we estimated, more
subjects will be needed; if the true effect is larger, fewer.  The
sub-sampling curves should therefore be read as ``what would happen
if the world is exactly as our 26 subjects suggest''.

%==============================================================================
\section{Synthesis}
%==============================================================================

Putting the eight sections together:

\begin{enumerate}
\item \textbf{Oberlin in the full-sentence block is the headline.}
      Accuracy is reliably above chance ($0.526$, $p=0.005$, $d=0.99$),
      and the result replicates across four
      alignment statistics ($d'$, $\cos(\mathrm{CI})$,
      $\beta_{\text{truetarget}}$, with the spectral cosine pointing
      the same way at $d=0.38$).  The convergence across statistics is good evidence that this is signal, not noise.

\item \textbf{UChicago in the full-sentence block is at chance.}
      Every metric is near zero with a small effect size.

\item \textbf{The site difference is trending in the predicted
      direction with a medium effect size} ($d=0.65$, $p=0.15$ in
      full-sentence) and is consistent with a real but under-powered
      site effect.  

\item \textbf{Full-sentence is numerically higher than imagined-sentence
      in every group} ($+0.014$ at UChicago, $+0.016$ at Oberlin,
      $+0.015$ Unified).  The within-subject difference does not
      reach $p < 0.05$ in any group, but the direction is consistent
      and the unified Cohen's $d$ is $+0.29$.  At $N=26$, $80\%$
      power for an effect of this size would require roughly $N=93$.

\item \textbf{Block order does not produce reliable effects.}  The
      strongest single trend is in Oberlin's imagined-sentence block
      ($p=0.08$ with $n_{B2}=5$), which is a practice
      pattern one might expect but is not significant.

\item \textbf{Questionnaire correlates are exploratory.}  No single
      $\rho$ survives Bonferroni.  BAIS\_C (auditory imagery
      \emph{control}) shows the cleanest cross-group descriptive
      pattern in the full-sentence block, and Book 3 of this project
      separately suggests BAIS\_V (auditory imagery
      \emph{vividness}) tracks individual significance.  Both point
      towards an auditory-imagery component to performance; neither
      is confirmed.
\end{enumerate}

\bigskip
\noindent

\begin{thebibliography}{1}
\bibitem{cohen1988} Cohen, J. (1988).
  \emph{Statistical power analysis for the behavioral sciences.}
  Lawrence Erlbaum.
\end{thebibliography}

\end{document}
